
\documentclass{article}
\usepackage{colortbl}
\usepackage{makecell}
\usepackage{multirow}
\usepackage{supertabular}

\begin{document}

\newcounter{utterance}

\centering \large Interaction Transcript for game `dond', experiment `semi', episode 3 with mistral{-}large{-}3{-}azure.
\vspace{24pt}

{ \footnotesize  \setcounter{utterance}{1}
\setlength{\tabcolsep}{0pt}
\begin{supertabular}{c@{$\;$}|p{.15\linewidth}@{}p{.15\linewidth}p{.15\linewidth}p{.15\linewidth}p{.15\linewidth}p{.15\linewidth}}
    \# & $\;$A & \multicolumn{4}{c}{Game Master} & $\;\:$B\\
    \hline

    \theutterance \stepcounter{utterance}  
    & & \multicolumn{4}{p{0.6\linewidth}}{
        \cellcolor[rgb]{0.9,0.9,0.9}{
            \makecell[{{p{\linewidth}}}]{
                \texttt{\tiny{[P1$\langle$GM]}}
                \texttt{Žaidžiate kooperatyvinį derybų žaidimą, kuriame turite sutarti su kitu žaidėju, kaip padalinti daiktų rinkinį tarp jūsų.} \\
\\ 
\texttt{Taisyklės:} \\
\texttt{(a) Jūs ir kitas žaidėjas gaunate daiktų rinkinį. Kiekvienam iš jūsų taip pat suteikiama slapta vertės funkcija, vaizduojanti, kiek vertinate kiekvieną daiktų tipą.} \\
\texttt{(b) Jūs keičiatės žinutėmis su kitu žaidėju, kad sutartumėte, kas gauna kokius daiktus. Kiekvienas gali išsiųsti ne daugiau nei 5 žinučių arba pabeigti žaidimą anksčiau, bet kuriuo metu pateikdamas slapią pasiūlymą.} \\
\texttt{(c) Jūs prašomi pateikti slaptą pasiūlymą, nurodantį norimus daiktus, suformatuotus skliaustuose taip: "{[}Pasiūlymas: <skaičius> <daikto pavadinimas>, <skaičius> <daikto pavadinimas>, <...>{]}"} \\
\texttt{(d) Jei jūsų pasiūlymai yra papildomi, t. y. yra pakankamai daiktų abiem pasiūlymams įvykdyti, kiekvienas žaidėjas gauna balus, remiantis gautų daiktų verčių suma. Priešingu atveju jūs abu gausite nulinį balą.} \\
\texttt{(e) Jūsų tikslas yra maksimizuoti savo balo ir kito žaidėjo balo sumą.} \\
\\ 
\texttt{Pradėkime.} \\
\\ 
\texttt{Prieinamų daiktų rinkinys yra:} \\
\\ 
\texttt{1 replės, 2 plytų, 1 žiedas, 1 laikrodis} \\
\\ 
\texttt{Jūsų slapta vertės funkcija:} \\
\\ 
\texttt{replės: 1, plytas: 0, žiedas: 2, laikrodis: 7} \\
\\ 
\texttt{SVARBU: Jūsų žinutės, išskyrus slaptą pasiūlymą, perduodamos tiesiogiai kitam žaidėjui, todėl neįtraukite jokio atsakymo į taisykles ar teksto, kuriame nurodomas jūsų pranešimas. Norėdami pateikti slaptą pasiūlymą, naudokite nurodytą formatą. Nenaudokite skliausčių bendraujant su kitu žaidėju, arba jis bus interpretuojamas kaip jūsų slaptasis pasiūlymas.} \\
            }
        }
    }
    & \\ \\

    \theutterance \stepcounter{utterance}  
    & \multicolumn{4}{p{0.6\linewidth}}{
        \cellcolor[rgb]{0.9,0.9,0.9}{
            \makecell[{{p{\linewidth}}}]{
                \texttt{\tiny{[P1$\rangle$GM]}}
                \texttt{Kokius daiktus vertinate labiausiai? Ar turite kokių nors prioritetų dalijantis šiuos daiktus?} \\
            }
        }
    }
    & & \\ \\

    \theutterance \stepcounter{utterance}  
    & & & \multicolumn{2}{p{0.3\linewidth}}{
        \cellcolor[rgb]{0.9,0.9,0.9}{
            \makecell[{{p{\linewidth}}}]{
                \texttt{\tiny{[GM$|$GM]}}
                \texttt{continue} \\
            }
        }
    }
    & & \\ \\

    \theutterance \stepcounter{utterance}  
    & & & \multicolumn{2}{p{0.3\linewidth}}{
        \cellcolor[rgb]{0.9,0.9,0.9}{
            \makecell[{{p{\linewidth}}}]{
                \texttt{\tiny{[GM$|$GM]}}
                \texttt{Kokius daiktus vertinate labiausiai? Ar turite kokių nors prioritetų dalijantis šiuos daiktus?} \\
            }
        }
    }
    & & \\ \\

    \theutterance \stepcounter{utterance}  
    & & \multicolumn{4}{p{0.6\linewidth}}{
        \cellcolor[rgb]{0.9,0.9,0.9}{
            \makecell[{{p{\linewidth}}}]{
                \texttt{\tiny{[GM$\rangle$P2]}}
                \texttt{Žaidžiate kooperatyvinį derybų žaidimą, kuriame turite sutarti su kitu žaidėju, kaip padalinti daiktų rinkinį tarp jūsų.} \\
\\ 
\texttt{Taisyklės:} \\
\texttt{(a) Jūs ir kitas žaidėjas gaunate daiktų rinkinį. Kiekvienam iš jūsų taip pat suteikiama slapta vertės funkcija, vaizduojanti, kiek vertinate kiekvieną daiktų tipą.} \\
\texttt{(b) Jūs keičiatės žinutėmis su kitu žaidėju, kad sutartumėte, kas gauna kokius daiktus. Kiekvienas gali išsiųsti ne daugiau nei 5 žinučių arba pabeigti žaidimą anksčiau, bet kuriuo metu pateikdamas slapią pasiūlymą.} \\
\texttt{(c) Jūs prašomi pateikti slaptą pasiūlymą, nurodantį norimus daiktus, suformatuotus skliaustuose taip: "{[}Pasiūlymas: <skaičius> <daikto pavadinimas>, <skaičius> <daikto pavadinimas>, <...>{]}"} \\
\texttt{(d) Jei jūsų pasiūlymai yra papildomi, t. y. yra pakankamai daiktų abiem pasiūlymams įvykdyti, kiekvienas žaidėjas gauna balus, remiantis gautų daiktų verčių suma. Priešingu atveju jūs abu gausite nulinį balą.} \\
\texttt{(e) Jūsų tikslas yra maksimizuoti savo balo ir kito žaidėjo balo sumą.} \\
\\ 
\texttt{Pradėkime.} \\
\\ 
\texttt{Prieinamų daiktų rinkinys yra:} \\
\\ 
\texttt{1 replės, 2 plytų, 1 žiedas, 1 laikrodis} \\
\\ 
\texttt{Jūsų slapta vertės funkcija:} \\
\\ 
\texttt{replės: 1, plytas: 1, žiedas: 4, laikrodis: 3} \\
\\ 
\texttt{SVARBU: Jūsų žinutės, išskyrus slaptą pasiūlymą, perduodamos tiesiogiai kitam žaidėjui, todėl neįtraukite jokio atsakymo į taisykles ar teksto, kuriame nurodomas jūsų pranešimas. Norėdami pateikti slaptą pasiūlymą, naudokite nurodytą formatą. Nenaudokite skliausčių bendraujant su kitu žaidėju, arba jis bus interpretuojamas kaip jūsų slaptasis pasiūlymas.} \\
\\ 
\texttt{Kokius daiktus vertinate labiausiai? Ar turite kokių nors prioritetų dalijantis šiuos daiktus?} \\
            }
        }
    }
    & \\ \\

    \theutterance \stepcounter{utterance}  
    & & & \multicolumn{4}{p{0.6\linewidth}}{
        \cellcolor[rgb]{0.9,0.9,0.9}{
            \makecell[{{p{\linewidth}}}]{
                \texttt{\tiny{[GM$\langle$P2]}}
                \texttt{Žiedas ir laikrodis yra vertingiausi man. Plytos ir replės yra mažiau vertingos. Kokie tavo prioritetai?} \\
            }
        }
    }
     \\ \\

    \theutterance \stepcounter{utterance}  
    & & & \multicolumn{2}{p{0.3\linewidth}}{
        \cellcolor[rgb]{0.9,0.9,0.9}{
            \makecell[{{p{\linewidth}}}]{
                \texttt{\tiny{[GM$|$GM]}}
                \texttt{continue} \\
            }
        }
    }
    & & \\ \\

    \theutterance \stepcounter{utterance}  
    & & & \multicolumn{2}{p{0.3\linewidth}}{
        \cellcolor[rgb]{0.9,0.9,0.9}{
            \makecell[{{p{\linewidth}}}]{
                \texttt{\tiny{[GM$|$GM]}}
                \texttt{Žiedas ir laikrodis yra vertingiausi man. Plytos ir replės yra mažiau vertingos. Kokie tavo prioritetai?} \\
            }
        }
    }
    & & \\ \\

    \theutterance \stepcounter{utterance}  
    & & \multicolumn{4}{p{0.6\linewidth}}{
        \cellcolor[rgb]{0.9,0.9,0.9}{
            \makecell[{{p{\linewidth}}}]{
                \texttt{\tiny{[P1$\langle$GM]}}
                \texttt{Žiedas ir laikrodis yra vertingiausi man. Plytos ir replės yra mažiau vertingos. Kokie tavo prioritetai?} \\
            }
        }
    }
    & \\ \\

    \theutterance \stepcounter{utterance}  
    & \multicolumn{4}{p{0.6\linewidth}}{
        \cellcolor[rgb]{0.9,0.9,0.9}{
            \makecell[{{p{\linewidth}}}]{
                \texttt{\tiny{[P1$\rangle$GM]}}
                \texttt{Aš labiausiai vertinu laikrodį ir žiedą, plytos man visiškai nevertingos. Ar norėtum pasidalinti žiedą ir laikrodį, o plytas ir replę palikti man?} \\
            }
        }
    }
    & & \\ \\

    \theutterance \stepcounter{utterance}  
    & & & \multicolumn{2}{p{0.3\linewidth}}{
        \cellcolor[rgb]{0.9,0.9,0.9}{
            \makecell[{{p{\linewidth}}}]{
                \texttt{\tiny{[GM$|$GM]}}
                \texttt{continue} \\
            }
        }
    }
    & & \\ \\

    \theutterance \stepcounter{utterance}  
    & & & \multicolumn{2}{p{0.3\linewidth}}{
        \cellcolor[rgb]{0.9,0.9,0.9}{
            \makecell[{{p{\linewidth}}}]{
                \texttt{\tiny{[GM$|$GM]}}
                \texttt{Aš labiausiai vertinu laikrodį ir žiedą, plytos man visiškai nevertingos. Ar norėtum pasidalinti žiedą ir laikrodį, o plytas ir replę palikti man?} \\
            }
        }
    }
    & & \\ \\

    \theutterance \stepcounter{utterance}  
    & & \multicolumn{4}{p{0.6\linewidth}}{
        \cellcolor[rgb]{0.9,0.9,0.9}{
            \makecell[{{p{\linewidth}}}]{
                \texttt{\tiny{[GM$\rangle$P2]}}
                \texttt{Aš labiausiai vertinu laikrodį ir žiedą, plytos man visiškai nevertingos. Ar norėtum pasidalinti žiedą ir laikrodį, o plytas ir replę palikti man?} \\
            }
        }
    }
    & \\ \\

    \theutterance \stepcounter{utterance}  
    & & & \multicolumn{4}{p{0.6\linewidth}}{
        \cellcolor[rgb]{0.9,0.9,0.9}{
            \makecell[{{p{\linewidth}}}]{
                \texttt{\tiny{[GM$\langle$P2]}}
                \texttt{Geriau laikrodį pasiimčiau aš, o žiedą atiduočiau tau. Plytas ir replę gali pasilikti. Ar tinka?} \\
            }
        }
    }
     \\ \\

    \theutterance \stepcounter{utterance}  
    & & & \multicolumn{2}{p{0.3\linewidth}}{
        \cellcolor[rgb]{0.9,0.9,0.9}{
            \makecell[{{p{\linewidth}}}]{
                \texttt{\tiny{[GM$|$GM]}}
                \texttt{continue} \\
            }
        }
    }
    & & \\ \\

    \theutterance \stepcounter{utterance}  
    & & & \multicolumn{2}{p{0.3\linewidth}}{
        \cellcolor[rgb]{0.9,0.9,0.9}{
            \makecell[{{p{\linewidth}}}]{
                \texttt{\tiny{[GM$|$GM]}}
                \texttt{Geriau laikrodį pasiimčiau aš, o žiedą atiduočiau tau. Plytas ir replę gali pasilikti. Ar tinka?} \\
            }
        }
    }
    & & \\ \\

    \theutterance \stepcounter{utterance}  
    & & \multicolumn{4}{p{0.6\linewidth}}{
        \cellcolor[rgb]{0.9,0.9,0.9}{
            \makecell[{{p{\linewidth}}}]{
                \texttt{\tiny{[P1$\langle$GM]}}
                \texttt{Geriau laikrodį pasiimčiau aš, o žiedą atiduočiau tau. Plytas ir replę gali pasilikti. Ar tinka?} \\
            }
        }
    }
    & \\ \\

    \theutterance \stepcounter{utterance}  
    & \multicolumn{4}{p{0.6\linewidth}}{
        \cellcolor[rgb]{0.9,0.9,0.9}{
            \makecell[{{p{\linewidth}}}]{
                \texttt{\tiny{[P1$\rangle$GM]}}
                \texttt{{[}Pasiūlymas: 1 laikrodis, 1 replė{]}} \\
            }
        }
    }
    & & \\ \\

    \theutterance \stepcounter{utterance}  
    & & & \multicolumn{2}{p{0.3\linewidth}}{
        \cellcolor[rgb]{0.9,0.9,0.9}{
            \makecell[{{p{\linewidth}}}]{
                \texttt{\tiny{[GM$|$GM]}}
                \texttt{proposal} \\
            }
        }
    }
    & & \\ \\

    \theutterance \stepcounter{utterance}  
    & & & \multicolumn{2}{p{0.3\linewidth}}{
        \cellcolor[rgb]{0.9,0.9,0.9}{
            \makecell[{{p{\linewidth}}}]{
                \texttt{\tiny{[GM$|$GM]}}
                [1, 0, 0, 1]
            }
        }
    }
    & & \\ \\

    \theutterance \stepcounter{utterance}  
    & & \multicolumn{4}{p{0.6\linewidth}}{
        \cellcolor[rgb]{0.9,0.9,0.9}{
            \makecell[{{p{\linewidth}}}]{
                \texttt{\tiny{[GM$\rangle$P2]}}
                \texttt{Kitas žaidėjas pateikė slaptą pasiūlymą. Dabar jūsų eilė pateikti savo pasiūlymą.} \\
            }
        }
    }
    & \\ \\

    \theutterance \stepcounter{utterance}  
    & & & \multicolumn{4}{p{0.6\linewidth}}{
        \cellcolor[rgb]{0.9,0.9,0.9}{
            \makecell[{{p{\linewidth}}}]{
                \texttt{\tiny{[GM$\langle$P2]}}
                \texttt{{[}Pasiūlymas: 1 laikrodis, 1 replė{]}} \\
            }
        }
    }
     \\ \\

    \theutterance \stepcounter{utterance}  
    & & & \multicolumn{2}{p{0.3\linewidth}}{
        \cellcolor[rgb]{0.9,0.9,0.9}{
            \makecell[{{p{\linewidth}}}]{
                \texttt{\tiny{[GM$|$GM]}}
                \texttt{proposal} \\
            }
        }
    }
    & & \\ \\

    \theutterance \stepcounter{utterance}  
    & & & \multicolumn{2}{p{0.3\linewidth}}{
        \cellcolor[rgb]{0.9,0.9,0.9}{
            \makecell[{{p{\linewidth}}}]{
                \texttt{\tiny{[GM$|$GM]}}
                [1, 0, 0, 1]
            }
        }
    }
    & & \\ \\

    \theutterance \stepcounter{utterance}  
    & & & \multicolumn{2}{p{0.3\linewidth}}{
        \cellcolor[rgb]{0.9,0.9,0.9}{
            \makecell[{{p{\linewidth}}}]{
                \texttt{\tiny{[GM$|$GM]}}
                [[1, 0, 0, 1], [1, 0, 0, 1]]
            }
        }
    }
    & & \\ \\

\end{supertabular}
}

\end{document}
